
\documentclass[theme=dnd]{rpgbook}

\graphicspath{{img/}}

\title{The Sunken Temple}
\subtitle{An example adventure for typesetters}
\author{A. W. Riter}
 
\cover{enraged-monster}
\RpgCoverTransparency{0.3}
\RpgPartTransparency{0.4}

%% this is the GM's copy, so show the secrets (dnd-handout.tex leaves them hidden)
\RpgSwitch{ShowSecrets}{on}

\begin{document}
	\maketitle

	%% D&D themed table of contents are possible, but with an example document this short it looks bad, so use it when it makes sense!
	% \tableofcontents


	\chapter{The Village of Saltmere}

		\begin{RpgNarration}
			The tide is out, and the whole village smells of salt and tar. Nets hang drying between the houses, and a bell clangs somewhere out on the flats, though there is no wind to ring it.

			The buildings around are dilapidated and empty, the only warmth coming from a low-slung building to the west, balanced on stilts which keep it out of the muddy water. The sign outside declares itself to be ``The Drowned Bell".
		\end{RpgNarration}
		%%%% This document is shared by both the dnd-book and the dnd-handout example



\RpgDropCap{S}{altmere} is a fishing village of forty-odd families, built on stilts above the mudflats where the river meets the sea. Its people are friendly enough to strangers who pay in coin, and wary of anyone who asks about the old temple.



\begin{RpgMap}[prefix={}]
	\area[pub]{The Drowned Bell}

	The village's only tavern stands at the end of the longest jetty. Its keeper, Morwen Tull, serves fish stew and weak ale, and knows everyone's business.

	\begin{RpgSecret}{Morwen is happy to talk about the temple, if the party buys a round.}
		Morwen's brother vanished into the temple last spring. She will offer the party 50 gold pieces to find out what happened to him. His body is found inside the temple, killed by the Typesetter.
	\end{RpgSecret}
	\area{The Dock} A grandiose name for the portion of the jetty facing out to sea, it is here that the old sea dogs gather before heading off to the pub for the evening, and trade the latest catch amongst themselves.

	\begin{RpgSecret*}
		Determine what the fishermen have caught using the following table:
		\begin{RpgTable}{llX}
			d4	& Fish & Result
			\\
			1 &		Nothing & The fishermen are angry. If the players enter the inn, \RpgDice{1d6+2} commoners will attempt to pick a fight with them.
			\\
			2 &		Minnows &  The mood in the inn is more sombre than usual.
			\\
			3 &		Cod & The fishermen are pleased. They buy all the players a round.
			\\
			4 &		Tuna & A raucous night is had by all. Any players who partake in the celebrations must succeed on a DC 13 Constitution saving throw or become \textit{poisoned} for one hour, but are immune to fear until the end of their next long rest.
		\end{RpgTable}

	\end{RpgSecret*}

	\begin{RpgArea}{The Chapel}
		A small chapel dedicated to the god Knonald Duth, run by the local priest, who is found within the \RpgMapRef{pub} (\RpgMapRef*{pub}) whenever not giving sermons.
		\begin{RpgSecret*}{}
			%% DM version
			\area{The Hidden Vestry} Behind a large portrait of Knonald is a door which leads to a dusty old room, where the priests would once have changed into ceremonial robes.
			\area{The Basement} In the centre of this cramped space is a book which has been bound to a pedestal and chained shut (DC 15 to unlock or force open). An individual who reads the book becomes \textit{incapacitated} until shaken awake by another creature.

			The book was written with Knonald's Quill.
		\end{RpgSecret*}

	\end{RpgArea}
	\area{The Sunken Temple} A forbidden temple which has fallen into disrepair and slipped beneath the waves. The townsfolk avoid it.

	\begin{RpgSecret*}
		\begin{RpgNarration}
			As the doors to the temple crack open, your eyes become fixed on the waves of beautiful writing which cover every surface of the temple's walls.

			It takes you a minute before you spot the emaciated bodies lining the walls, those unlucky enough to have become entranced by the writing, and stayed there until they starved to death, or worse, were killed by the hulking monster that now approaches you from the bowels of the temple...
		\end{RpgNarration}

		The Typesetter approaches! Roll initiative!
	\end{RpgSecret*}
\end{RpgMap}

		\appendix

	\chapter{Creatures}

		%% The adventure's final foe: a full-width (RpgStat*) legendary statblock
		\begin{RpgStat*}{The Typesetter}[
				size=Medium,
				type=Undead,
				alignment={Chaotic evil},
				AC=15,
				HP=12d8+24,
				str=12, dex=15, con=14, int=18, wis=14, cha=12,
				str-save, int-save,
				skills={Perception, Arcana},
				skills-expertise=History,
				cr=6,
				languages={Common, Draconic, \LaTeX},
				filigree,
			]
			\RpgStatLegendaryResistance{2}

			\begin{RpgStatSpells}[title={Drowned Spellcasting}, innate, modifier=int]
				\atWill{faulty preamble, thunderclap}
				\perDay{2}{fireball}
				\level[3][2]{compile document}
			\end{RpgStatSpells}

			\section{Actions}
				\attack{Inkblade}[distance=both, reach=10, range=20/60, modifier=int, dmg=1d8+\RpgStatMod{int}, dmg-type=necrotic]

				\RpgStatLegendarySection[uses=2]
				\action{Rapid Edit} \name{} makes an Inkblade attack, then moves up to 30 feet.

				\RpgStatMythicSection[narration={Just as the Typesetter falls, a voice rings out: \textit{``but will it compile on Overleaf?''} The ink on the walls begins to move\ldots}]
				\save{Online Woes}[modifier=int, dc=15, shape=sphere, origin={within 10 feet of the Typesetter}, success={half damage, and is not stunned}, fail={\RpgDice{4d10} psychic damage, and is stunned until the end of its next turn}]
		\end{RpgStat*}

		\begin{RpgStat}{Powerful Dragon Format}[
				size=Tiny,
				type=Dragon,
				alignment={Lawful neutral},
				AC=13,
				HP=2d8+2,
				speed={10ft, fly 30ft},
				str=5, dex=16, con=12, int=16, wis=12, cha=15,
				cr=1/8,
				int-save, wis-save,
				languages={Draconic, \LaTeX},
				nickname=PDF,
			]
			\trait{Portable} The \name{} is unaffected by its surroundings, ignoring difficult terrain and strong winds.

			\section{Actions}
				\attack{Claw}[modifier=dex, dmg-type=slashing, plus-dmg=1d6, plus-type=fire]
		\end{RpgStat}

		Characters who search the scriptorium after the Typesetter's defeat may roll once on the table below.

		%% A filigree frame can be put around anything: here, a loot table
		\begin{RpgFiligree}
			\begin{RpgTable}[title={The Typesetter's Hoard}]{lX}
				d6 & Loot
				\\ 1 & A case of lead type, worth \RpgDice{3d6} gold pieces to any printer.
				\\ 2 & A bottle of ink that never dries, and never fades.
				\\ 3 & A ream of vellum, each sheet watermarked with a sea serpent.
				\\ 4 & A brass composing stick, engraved with the name \textit{Knonald}.
				\\ 5 & A sealed letter from Morwen's brother, addressed to his sister.
				\\ 6 & \textit{Knonald's Quill}, if the Typesetter did not use it in the fight.
			\end{RpgTable}
		\end{RpgFiligree}

	\chapter{Treasure \& Magic}

		The following items can be found in the temple, and the Typesetter casts the spells below.

		\section{Magic Items}

			\begin{RpgItem}{Knonald's Quill}[
					type={Wondrous item},
					rarity=Rare,
					requires-attunement={by a creature with an Intelligence of 13 or higher},
				]
				As an action, you may use this quill to write on any surface. The symbols you write are so beautifully typeset that a creature which reads them must succeed on a DC 13 Intelligence saving throw or be charmed by you, as if by the \textit{charm person} spell.
			\end{RpgItem}

		\section{Spells}

			\begin{RpgSpell}{Compile Document}[
					level=3,
					school=Conjuration,
					range={60 feet},
					area={30-foot cube},
					components=VS,
				]
				Roll \RpgDice{1d4}. That many documents within the area are converted into a \textit{Powerful Dragon Format}, which acts on its own initiative and regards you as an ally. For each dragon created in this way, you regain 5 hit points.

				\RpgSpellUpcast[The hit points restored by each dragon increase by]{5}
			\end{RpgSpell}
			\begin{RpgSpell}{Faulty Preamble}[
					school=Evocation,
					range=Touch,
					duration={Until dispelled},
				]
				You corrupt a document that you touch. When this document is next affected by \textit{compile document}, it does not become a Powerful Dragon Format. Instead, it deals \RpgDice{1d10} psychic damage to the creature attempting the compilation.

				\RpgCantripScaling{1d10}{2d10}{3d10}{4d10}
			\end{RpgSpell}
\end{document}
